\documentclass{article}
\usepackage[a4paper, margin=2cm]{geometry} 
\usepackage[T1]{fontenc} %lädt Deutsche schriftzeichen
\usepackage{amssymb}%für Mathe
\usepackage{setspace} %für Zeilen Abstand 
\usepackage[none]{hyphenat} % Silbentrennung deaktivieren 
\usepackage{fancyhdr}%Kopf und Fußzeile: https://de.overleaf.com/learn/latex/Headers_and_footers
\usepackage{graphicx}%Zum Bilder einfügen
\usepackage{xcolor}
\usepackage{wrapfig}    
\usepackage[ddmmyyyy]{datetime}%zum bearbeiten vom Datungs Format
\usepackage{eurosym}%für \Euro / \euro
\usepackage{hyperref}



\newcommand{\Thema}{Formelsammlung}
\newcommand{\Fach}{Physik}
\newcommand{\Zeilenabstand}{1.2}
\newcommand{\Inhaltsverzeichnisname}{Inhaltsverzeichnis}

\newcommand{\absatz}{5mm}
\newcommand{\Blocksatz}{\justifying}
\newcommand{\strich}{\vspace{\absatz}\hrule\vspace{\absatz}}


%For subsubsub... sections
\setcounter{secnumdepth}{5}
\setcounter{tocdepth}{5}

\makeatletter
\newcommand\subsubsubsection{\@startsection{paragraph}{4}{\z@}{-2.5ex\@plus -1ex \@minus -.25ex}{1.25ex \@plus .25ex}{\normalfont\normalsize\bfseries}}
\newcommand\subsubsubsubsection{\@startsection{subparagraph}{5}{\z@}{-2.5ex\@plus -1ex \@minus -.25ex}{1.25ex \@plus .25ex}{\normalfont\normalsize\bfseries}}
\makeatother


\title{\Thema}
\author{Niklas Hainz}
\date{\today}

\setlength{\parindent}{0cm} %deaktiviert die Automatische Einrückung am Anfang des Dokumentes 
\renewcommand{\contentsname}{\Inhaltsverzeichnisname}







\begin{document}
    \begin{spacing}{\Zeilenabstand}
        \maketitle
        \begin{center}
            \vspace{-5mm}
            Fach: \Fach
        \end{center}




        \newpage

        \pagestyle{fancy}
        \fancyhead[HL]{Niklas Hainz}
        \fancyhead[HC]{\Thema / Fach: \Fach}
        \fancyhead[HR]{\today}

        \tableofcontents
        \newpage

        \section{SI-Basisgrößen und -einheiten}


            \begin{center}
                \begin{table}[h]
                    \begin{tabular}{|l|l|l|l|}
                        \multicolumn{2}{l}{\textbf{Basisgröße}}     & \multicolumn{2}{l}{\textbf{Basiseinheit}} \\
                        \hline
                        \textbf{Name}           & \textbf{Symbol}   & \textbf{Name}      & \textbf{Symbol}      \\
                        \hline
                        Zeit                    & t                 & Sekunde            & s                    \\
                        \hline
                        Länge                   & l                 & Meter              & m                    \\
                        \hline
                        Masse                   & m                 & Kilogramm          & kg                   \\
                        \hline
                        elektrische Stromstärke & I                 & Ampere             & A                    \\
                        \hline
                        Temperatur              & T                 & Kelvin             & K                    \\
                        \hline
                        Stoffmenge              & n                 & Mol                & mol                  \\
                        \hline
                        Lichtstärke             & I\textsubscript{v}& Candela            & cd                   \\
                        \hline
                    \end{tabular}
                \end{table}
            \end{center}


        \section{Mechanik}
            \subsection{Gleichförmige Bewegung}  
                Formel für Gleichförmige Bewegungen, diese Werden benutz wenn die Geschwindigkeit(\(v\)) \textbf{konstant} ist bzw. 
                wenn keine beschleunigung in der Bewegung vorhanden ist.  

                \begin{equation}
                    s=v\cdot t
                \end{equation}

            \subsection{Gleichmäßig beschleunigte Bewegung}
                Wird genuzt wenn die Beschleunigung(\(a\)) von der Bewegung konstant ist 

                \begin{equation}
                    s=\frac{1}{2}\cdot a \cdot t^2
                \end{equation}

                \begin{equation}
                    s=\frac{1}{2}\cdot v \cdot t 
                \end{equation}

            \subsection{Kräfte}

            \subsubsection{Kräfte Berechnung}

            Um die Kraft in Newton(N) zu berrechnen
            \begin{equation}
                F=m\cdot a
            \end{equation}

            \subsubsection{Defintion Newton}

            Die Definition von der Einheit Newton(\(N\)) ist:

            \begin{equation}
                N=kg\cdot \frac{m}{s^{2}}
            \end{equation}

            \subsection{Schiefe Ebene}




            \newpage
            \subsection{Arbeit, Energie, Leisung}

            \subsubsection{Einheiten}

                Arbeit W (Work)
            
                Energie E (Energie) \hfil [W] (Watt) = [E] = J (Joule) = Ws (Wattsekunde) = Nm (Newtonmeter)
            
                Leisung $P=\frac{W}{T}$  
            
                [P] = w $\Rightarrow$ "W" als \textbf{Einheit} ist watt. Als formelzeichen Work 

            \subsubsection{Formeln}

            \begin{center}
            \begin{tabular}{|l | l| l |}
                \hline
                 & Mechanische Arbeit & $W=F\cdot s$ \\
                \hline
                Arbeit W  & Hubarbeit & $W_H=mgh$ \\
                \hline
                & beschleunigungsarbeit & $W_B=mas$  \\
                \hline
            \end{tabular}

            \strich

            \begin{tabular}{|l| l| l|}
                \hline
                & Lageenergie & $E_{pot}=mgh$ \\
                \hline
                Energie E & Bewegungsenergie & $E_{kin} = \frac{1}{2}mv^2 $\\
                \hline
                & Elektrische Energie & $E_{el}=PT=UIT$\\
                \hline
            \end{tabular}

            \strich

            \begin{tabular}{|l| l| l|}
                \hline
                & Mechanische Leistung & $P_{mesch}=\frac{Fs}{t}$\\
                \hline
                Leistung $P=\frac{W}{t}$ & Elektrische Leistung & $P_{el}=UI$ \\
                \hline

            \end{tabular}

            \end{center}


            \subsection{Wirkungsgrad}

                Um den Wirkungsgrad zu berrechnen:
                \begin{equation}
                    \eta_{ges} = \eta_1 \cdot \eta_2 \ldots
                \end{equation}



            \subsection{Wagerrechter Wurf}

                \url{https://www.leifiphysik.de/mechanik/waagerechter-und-schraeger-wurf/grundwissen/waagerechter-wurf}


            \subsection{Schräger Wurf}

                \url{https://www.leifiphysik.de/mechanik/waagerechter-und-schraeger-wurf/grundwissen/schraeger-wurf-nach-oben-mit-anfangshoehew}


        \section{Elektrisches und magnetisches Feld}

            \subsection{Elektrisches Feld}

                \begin{equation}
                    1C = [A_{s}]
                \end{equation}

                \begin{equation}
                    F_{el} = F_{c} = \frac{1}{4\pi \epsilon} \cdot \frac{q_{1}\cdot q_{2}}{r^{2}}
                \end{equation}
                \begin{equation}
                    E=\frac{F_{el}}{q}
                \end{equation}
                \begin{equation}
                    E= \frac{1}{4\pi\epsilon}\cdot \frac{Q}{r^{2}}
                \end{equation}

                \subsubsection{Monotones Feld}

                    \url{https://www.leifiphysik.de/elektrizitaetslehre/ladungen-elektrisches-feld/grundwissen/homogenes-elektrisches-feld}

				\subsubsection{Kondensator \& Kapazität, Homogenes Feld}

					\begin{equation}
						C = \frac{Q}{U} = \epsilon_{0}  \cdot \epsilon_{r} \frac{A}{d}
					\end{equation}
					% Formel für die kapazität C im bezug auf den Plattenkondensator
					\begin{equation}
						E = \frac{U}{d}
					\end{equation}
					\hspace{15mm} El. Energie
					\begin{equation}
						W_{el} = \frac{1}{2} \cdot C \cdot U^2
					\end{equation}
					\hspace{15mm} El. Feld
					\begin{equation}
						E = \frac{\frac{Q}{C}}{d} = \frac{Q}{Cd}
					\end{equation}
					



    \end{spacing}
\end{document}